% Topic 1.2.05 — Limit Theorems (Chebyshev, LLN, CLT).

\section{Limit Theorems (Chebyshev, LLN, CLT)}
\label{sec:1.2.05-limit-theorems}

\begin{ticket}
Fundamental theorems of probability. Chebyshev's inequality. Law of Large Numbers. Central Limit Theorem.
\end{ticket}

\subsection*{Theory}

We turn from properties of fixed distributions to the behaviour of
sums of many random variables. Three pillars structure the
discussion: tail inequalities (Markov, Chebyshev) that bound rare
deviations; the law of large numbers, which says that empirical
averages concentrate around the mean; and the central limit theorem,
which says that the fluctuations around that mean are universally
Gaussian after rescaling.

\begin{theorem}[Markov's inequality]
\label{thm:markov-ineq}
For any non-negative random variable $X \geq 0$ and any $a > 0$,
\[
  \Prob(X \geq a) \;\leq\; \frac{\E X}{a}.
\]
\end{theorem}
\proofof{markov-ineq}

\begin{theorem}[Chebyshev's inequality]
\label{thm:chebyshev-ineq}
For any $X \in L^2$ and any $a > 0$,
\[
  \Prob\bigl(|X - \E X| \geq a\bigr) \;\leq\; \frac{\Var X}{a^2}.
\]
\end{theorem}
\proofof{chebyshev-ineq}

These are the workhorse tail bounds: distribution-free, requiring
only one or two moments. The constants are tight in the sense that
no sharper bound holds for the whole class of $L^1$ (resp.~$L^2$)
random variables — see \cite[Ch.~II, \S~1]{Shiryaev1989}.

\medskip
\noindent\textbf{Modes of convergence of random variables.} A
sequence $(Y_n)$ may approach a limit~$Y$ in several inequivalent
senses.

\begin{definition}[Modes of convergence]
\label{def:convergence-modes}
Let $Y_1, Y_2, \dots, Y$ be random variables on a common probability
space.
\begin{itemize}
  \item \emph{In probability,} $Y_n \xrightarrow{\Prob} Y$, if for every
        $\varepsilon > 0$, $\Prob(|Y_n - Y| > \varepsilon) \to 0$.
  \item \emph{Almost surely,} $Y_n \xrightarrow{\text{a.s.}} Y$, if
        $\Prob(\{\omega : Y_n(\omega) \to Y(\omega)\}) = 1$.
  \item \emph{In distribution,} $Y_n \xrightarrow{d} Y$, if
        $F_{Y_n}(x) \to F_Y(x)$ at every point~$x$ of continuity
        of~$F_Y$.
\end{itemize}
\end{definition}

The implications are
$\xrightarrow{\text{a.s.}} \Rightarrow \xrightarrow{\Prob} \Rightarrow \xrightarrow{d}$;
none of the reverse implications hold in general. Convergence in
distribution depends only on the marginal laws and is the weakest
notion — it is the one used by the central limit theorem.

\begin{theorem}[Weak Law of Large Numbers]
\label{thm:wlln}
Let $X_1, X_2, \dots$ be independent random variables with common
mean $\E X_i = \mu$ and uniformly bounded variance $\Var X_i \leq C < \infty$.
Set $\bar X_n = (X_1 + \cdots + X_n)/n$. Then
\[
  \bar X_n \xrightarrow{\Prob} \mu \qquad (n \to \infty).
\]
In particular, when the $X_i$ are i.i.d.\ with finite variance, the
sample mean converges in probability to the common mean.
\end{theorem}
\proofof{wlln}

\begin{theorem}[Strong Law of Large Numbers, Kolmogorov]
\label{thm:slln}
Let $X_1, X_2, \dots$ be i.i.d.\ random variables with $\E |X_1| < \infty$
and $\E X_1 = \mu$. Then
\[
  \bar X_n \xrightarrow{\text{a.s.}} \mu \qquad (n \to \infty).
\]
\end{theorem}
\proofof{slln}

The strong law strengthens the weak law in two respects: (a) it
upgrades convergence in probability to almost-sure convergence, and
(b) it requires only a first moment instead of a second. The proof
is harder; we sketch only the case $\E X_1^4 < \infty$ via the
Borel–Cantelli lemma. For the general $L^1$ case, see
\cite[Ch.~III, \S~3]{Shiryaev1989} or \cite[Ch.~VI]{Gnedenko2001}.

\begin{theorem}[Central Limit Theorem, Lindeberg–Lévy]
\label{thm:clt}
Let $X_1, X_2, \dots$ be i.i.d.\ random variables with $\E X_1 = \mu$
and $\Var X_1 = \sigma^2 \in (0, \infty)$. Set $S_n = X_1 + \cdots + X_n$
and
\[
  Z_n \;=\; \frac{S_n - n\mu}{\sigma \sqrt{n}}.
\]
Then $Z_n \xrightarrow{d} \mathcal{N}(0, 1)$, i.e.\ for every
$x \in \R$,
\[
  \Prob(Z_n \leq x) \;\to\; \Phi(x) \;:=\; \frac{1}{\sqrt{2\pi}} \int_{-\infty}^{x} e^{-t^2/2}\, dt.
\]
\end{theorem}
\proofof{clt}

\begin{remark}
The strong law is asymptotic ($\bar X_n \to \mu$); Chebyshev is a
quantitative bound ($\Prob(|\bar X_n - \mu| \geq \varepsilon) \leq C/(n\varepsilon^2)$);
the CLT identifies the \emph{rate} of fluctuation as $1/\sqrt{n}$
with a \emph{universal} Gaussian limit. The clean proof of the CLT
via characteristic functions is taken up in
\cref{sec:1.2.06-standard-distributions}; here we only sketch the
key idea so the statement is in hand.
\end{remark}

\subsection*{Proofs}

\begin{proof}[\Cref{thm:markov-ineq}]
\proofanchor{markov-ineq}
Decompose $X = X \cdot \mathbf{1}_{\{X \geq a\}} + X \cdot \mathbf{1}_{\{X < a\}}$.
Both summands are non-negative (since $X \geq 0$), so by monotonicity
and linearity (\cref{prop:exp-properties}),
\[
  \E X \;\geq\; \E\bigl[X \cdot \mathbf{1}_{\{X \geq a\}}\bigr] \;\geq\; \E\bigl[a \cdot \mathbf{1}_{\{X \geq a\}}\bigr] \;=\; a\, \Prob(X \geq a).
\]
Divide by $a > 0$.
\end{proof}

\begin{proof}[\Cref{thm:chebyshev-ineq}]
\proofanchor{chebyshev-ineq}
Apply \cref{thm:markov-ineq} to the non-negative variable $(X - \E X)^2$
with threshold $a^2$:
\[
  \Prob\bigl(|X - \E X| \geq a\bigr) = \Prob\bigl((X - \E X)^2 \geq a^2\bigr) \leq \frac{\E[(X - \E X)^2]}{a^2} = \frac{\Var X}{a^2}. \qedhere
\]
\end{proof}

\begin{proof}[\Cref{thm:wlln}]
\proofanchor{wlln}
By independence and \cref{cor:var-sum-indep},
$\Var(X_1 + \cdots + X_n) = \sum_i \Var X_i \leq nC$, hence
$\Var \bar X_n = \Var(X_1 + \cdots + X_n)/n^2 \leq C/n$. By
\cref{thm:chebyshev-ineq} applied to $\bar X_n$ (which has mean
$\mu$ by linearity),
\[
  \Prob\bigl(|\bar X_n - \mu| \geq \varepsilon\bigr) \;\leq\; \frac{\Var \bar X_n}{\varepsilon^2} \;\leq\; \frac{C}{n \varepsilon^2} \;\to\; 0.
\]
\end{proof}

\begin{proof}[\Cref{thm:slln}, sketch under $\E X_1^4 < \infty$]
\proofanchor{slln}
Without loss of generality $\mu = 0$ (replace $X_i$ by $X_i - \mu$).
Set $S_n = X_1 + \cdots + X_n$. Expanding $\E S_n^4$ and using
i.i.d.\ with zero mean, only the $\E X_i^4$ and the
$\binom{n}{2} \binom{4}{2} \E[X_i^2] \E[X_j^2] = 3 n(n-1) \sigma^4$
terms survive, giving $\E S_n^4 \leq c\, n^2$ for a constant
$c$ depending on the fourth moment. By Markov,
$\Prob(|\bar X_n| \geq \varepsilon) = \Prob(|S_n|^4 \geq n^4 \varepsilon^4) \leq c\, n^2 / (n^4 \varepsilon^4) = c/(n^2 \varepsilon^4)$.
This is summable in $n$, so by the first Borel–Cantelli lemma
(\cite[Ch.~II, \S~10]{Shiryaev1989}) $\Prob(|\bar X_n| \geq \varepsilon \text{ infinitely often}) = 0$
for each $\varepsilon$, and intersecting over a sequence
$\varepsilon \to 0$ gives $\bar X_n \to 0$ a.s. The general
$L^1$ case requires Kolmogorov's truncation argument and is in
the references cited above.
\end{proof}

\begin{proof}[\Cref{thm:clt}, idea]
\proofanchor{clt}
Standardise: with $Y_i = (X_i - \mu)/\sigma$, the variables $Y_i$ are
i.i.d.\ with $\E Y_i = 0$, $\Var Y_i = 1$, and
$Z_n = (Y_1 + \cdots + Y_n)/\sqrt{n}$.

Let $\varphi(t) = \E e^{itY_1}$ be the characteristic function of
$Y_1$ (\cref{def:char-function}). Then
\[
  \E e^{itZ_n} \;=\; \E\, e^{it(Y_1 + \cdots + Y_n)/\sqrt n} \;=\; \varphi\bigl(t/\sqrt n\bigr)^n,
\]
using independence and the multiplicativity of $\E e^{i t \cdot \,}$
on independent sums (\cref{thm:indep-multiplicative}). The Taylor
expansion of $\varphi$ at the origin reads
$\varphi(s) = 1 - \tfrac{1}{2} s^2 + o(s^2)$ as $s \to 0$ (using
$\varphi'(0) = i \E Y_1 = 0$, $\varphi''(0) = -\E Y_1^2 = -1$).
Therefore
\[
  \varphi(t/\sqrt n)^n = \Bigl(1 - \tfrac{t^2}{2n} + o(1/n)\Bigr)^n \to e^{-t^2/2} \qquad (n \to \infty),
\]
which is the characteristic function of $\mathcal{N}(0, 1)$. By the
continuity theorem of Lévy (\cref{thm:levy-continuity}), pointwise
convergence of characteristic functions to a continuous limit
implies convergence in distribution to the corresponding law,
i.e.\ $Z_n \xrightarrow{d} \mathcal{N}(0, 1)$.
\end{proof}

\subsection*{Examples}

\begin{example}[Sample size from Chebyshev]
\label{ex:cheb-poll}
A pollster wants to estimate a proportion $p \in [0, 1]$ from
$n$ independent Bernoulli($p$) responses, with the empirical
proportion $\bar X_n$ within $\varepsilon = 0.03$ of~$p$ with
probability at least~$0.95$. Since
$\Var X_i = p(1-p) \leq 1/4$, Chebyshev gives
\[
  \Prob\bigl(|\bar X_n - p| \geq \varepsilon\bigr) \leq \frac{p(1-p)}{n \varepsilon^2} \leq \frac{1}{4 n \varepsilon^2}.
\]
For this to be $\leq 0.05$, take $n \geq 1 / (4 \cdot 0.05 \cdot \varepsilon^2) = 1/(0.2 \cdot 9 \cdot 10^{-4}) \approx 5556$.
The CLT-based bound (\cref{ex:clt-poll}) sharpens this dramatically.
\end{example}

\begin{example}[Empirical mean of fair coin flips, WLLN]
\label{ex:wlln-coin}
Let the $X_i$ be i.i.d.\ Bernoulli($1/2$). Then $\E X_i = 1/2$,
$\Var X_i = 1/4$, and \cref{thm:wlln} gives $\bar X_n \to 1/2$ in
probability with the rate
\[
  \Prob\bigl(|\bar X_n - 1/2| \geq 0.01\bigr) \leq \frac{1}{4 n \cdot 10^{-4}} = \frac{2500}{n};
\]
so for $n = 10^6$ the empirical proportion is within $0.01$ of
$1/2$ with probability at least $0.9975$. The strong law
(\cref{thm:slln}) further says that for almost every infinite
sequence of coin flips, the empirical proportion actually
\emph{converges} to $1/2$ — not merely concentrates around it.
\end{example}

\begin{example}[Normal approximation to a binomial, CLT]
\label{ex:clt-poll}
With the polling setup of \cref{ex:cheb-poll} and $p$ unknown but
$p(1-p) \leq 1/4$, the CLT gives
$Z_n = \sqrt n (\bar X_n - p)/\sqrt{p(1-p)} \approx \mathcal{N}(0,1)$,
so
\[
  \Prob\bigl(|\bar X_n - p| \geq \varepsilon\bigr) \approx 2\bigl(1 - \Phi(\varepsilon \sqrt n / \sqrt{p(1-p)})\bigr) \leq 2(1 - \Phi(2 \varepsilon \sqrt n)).
\]
Setting this $\leq 0.05$ requires $\Phi(2 \varepsilon \sqrt n) \geq 0.975$,
i.e.\ $2 \varepsilon \sqrt n \geq 1.96$, so $n \geq (1.96/(2 \varepsilon))^2 = (1.96/0.06)^2 \approx 1067$ —
about five times smaller than the Chebyshev sample size, and the
standard rule of thumb behind ``$\pm 3\%$, 95\% confidence, $n \approx 1100$''
in classical opinion polling.
\end{example}

\begin{problem}
\label{prob:dice-clt}
Roll a fair die $n = 100$ times and let $S$ be the sum of the
outcomes. Use the CLT to estimate $\Prob(S \geq 380)$.
\end{problem}

\begin{proof}[Solution]
Each outcome $X_i$ has $\E X_i = 7/2$ and $\Var X_i = 35/12$
(\cref{ex:die-moments}). With $n = 100$, $\E S = 350$ and
$\Var S = 100 \cdot 35/12 = 875/3$, so $\sigma_S = \sqrt{875/3} \approx 17.08$.
The standardised threshold is
\[
  z = \frac{380 - 350}{17.08} \approx 1.756,
\]
and by \cref{thm:clt}, $\Prob(S \geq 380) \approx 1 - \Phi(1.756) \approx 1 - 0.9605 \approx 0.0395$.
A continuity correction (using $379.5$ in place of $380$) gives
$z \approx 1.726$ and the slightly larger $\approx 0.0422$, in
better agreement with the exact (computer-evaluated) probability.
\end{proof}
